\section{Introduction}

Minimizing redundant computation is a fundamental concern in query optimization.
Query containment is a classic tool for this purpose: it determines whether the results of one query are always included in those of another, independently of the targeted database.
Solving this problem underpins a wide range of optimization techniques, including caching, view selection, query rewriting, and equivalence validation~\cite{Afrati2019-ch1, DOAN201221, Chaudhuri1993}.
Under set semantics, multiple query fragments have been thoroughly characterized~\cite{Chandra1977, aho1979equivalences, aho1979efficient, Chekuri2000}.
Bag-semantic and bag-set semantic containment have also been studied~\cite{Chaudhuri1993, Afrati2010, khamis2021bag}; however, several query fragments remain undecidable under these semantics and open problems persist.

\ac{KG} has become a widely adopted paradigm for representing and integrating heterogeneous data, enabling applications in scientific data management, enterprise knowledge representation, and the Web of Data.
SPARQL is the standard query language for \ac{KG} built on \ac{RDF}, and federated SPARQL extends this by enabling virtual data integration across distributed \ac{KG} sources, allowing heterogeneous data to be queried jointly without centralization, which is particularly valuable when sources are governed by different policies, maintained by independent parties, or too large to consolidate.
However, federated queries impose significant computational demands on both endpoints and clients, a challenge further compounded by the reliability limitations of public SPARQL endpoints, which frequently experience downtime or performance degradation under load~\cite{buil2013sparql}.
In settings such as scientific reproducibility studies, where queries must be re-executed or iteratively refined, this overhead is especially costly.
Query containment is therefore a natural optimization: when a new query is contained in a previously executed one, results can be derived from cached data without re-querying the federation, reducing both endpoint load and network traffic.

For SPARQL, existing research has addressed set-semantic query containment~\cite{WudageChekol2013}.
Yet the SPARQL query language operates under bag-set semantics~\cite{Perez2006, Angles2016}, in which the underlying database operates under set semantics~\cite{w3ConceptsAbstract} while query operators produce bags\footnote{A bag, or multiset, is similar to set but allows duplicate elements.}~\cite{w3SPARQLQuery}.
Set-semantic containment does not imply bag-set semantic containment, making existing containment results invalid for real-world SPARQL query optimization.

To the best of our knowledge, bag-set semantic SPARQL query containment remains unexplored, as does query containment in the context of federated queries.\footnote{Note that QCan~\cite{salas2018canonicalisation} deals with bag-set semantics, however, it addresses query equivalence rather than query containment, by transforming SPARQL queries into a canonical form and performing syntactic comparison.}

This paper bridges these gaps by formalizing query containment for conjunctive SPARQL federated queries under bag-set semantics.
We first present the \hyperref[sec:related_works]{related works},
then the \hyperref[sec:preliminaries]{preliminaries},
the \hyperref[sec:semantic_federated_q]{semantics of federated queries},
followed by the evaluation of \hyperref[sec:federated_queries_query_containment]{federated queries query containment}
and we \hyperref[sec:conclusion]{conclude} the paper.
